\chapter{電気的機能によるニューロンの分類}
\chaptermark{素子としてのニューロン}
\label{sec:eetypes}
\begin{chaptergoals}
\item ニューロンがどんな素子として働くか：積分器、フィルタ、変換器、発振器、ラッチ；
\item 1個の細胞がバンドパス共振器や、\nt{SERO}がオンにするラッチになる仕組み；
\item 漏れのあるニューロンからフィードバックで作る漏れのない積分器と、精密な調整が要る理由；
\item なぜこれらは型でなくモードであり、イオンチャネルとブロードキャストチャネルが選ぶのか。
\end{chaptergoals}

第\ref{sec:neurontypes}章では、出力が放出する神経伝達物質によってニューロンを分類した。電子技術者なら別の分け方をするだろう。ニューロンが自分の信号に何をするか、つまりどんな素子として働くかによる分類である。本書の主要な素子は第\ref{sec:glutcomputer}章から知っている。漏れ積分ニューロン~\eqref{eq:glutneuron}、すなわちコンパレータ付きのRC回路である。しかし脳には他の素子もある。表~\ref{tab:eetypes}にそれらを四つのグループにまとめた。ほとんどはすでに本書に登場している。

\begin{center}
\footnotesize
\setlength{\tabcolsep}{3pt}
\renewcommand{\arraystretch}{1.15}
\begin{longtable}{>{\raggedright}p{3.1cm}>{\raggedright}p{3.3cm}>{\raggedright}p{4.7cm}>{\raggedright\arraybackslash}p{2.4cm}}
\caption{素子としてのニューロン：名前、電子回路での対応物、素子の働き、本書での登場箇所。}\label{tab:eetypes}\\
\toprule
ニューロン & 素子 & 働き & 本書での箇所 \\
\midrule
\endfirsthead
\toprule
ニューロン & 素子 & 働き & 本書での箇所 \\
\midrule
\endhead
\bottomrule
\endfoot
\multicolumn{4}{l}{\emph{フィルタと積分器}} \\
漏れ積分ニューロン & コンパレータ付きRC積分器 & 窓$\tau$内の入力を足し合わせ、閾値で発火する & 第\ref{sec:glutcomputer}章、\eqref{eq:glutneuron} \\
同時検出器 & 時間窓付きAND & 入力がほぼ同時に届いたときだけ発火する。$\tau$がごく小さい & 第\ref{sec:glutcomputer}、\ref{sec:code}章、\eqref{eq:window} \\
相動性ニューロン & ハイパスフィルタ、エッジ検出器 & 入力の変化に応答し、その後沈黙する & 第\ref{sec:sensors}章 \\
順応するニューロン & 自動ゲイン制御付きハイパスフィルタ & 一定の入力では発火率が下がる & 第\ref{sec:sensors}章、\eqref{eq:nakarushton} \\
共振器 & 共振回路、バンドパスフィルタ & 自分の周波数の入力に最もよく応答する & 第\ref{sec:resonator}節 \\
漏れのない積分器 & オペアンプ積分器 & 速度を位置に変えて保持する & 第\ref{sec:perfectint}節 \\
\midrule
\multicolumn{4}{l}{\emph{変換器}} \\
持続性ニューロン & 電圧--周波数変換器 & 発火率が入力に線形に追従し、ほとんど疲れない & 第\ref{sec:code}章 \\
段階的電位ニューロン（スパイクなし） & アナログ増幅器、バッファ & 滑らかな電圧を短い距離だけ伝える & 第\ref{sec:sensors}章 \\
整流ニューロン & 半波整流器 & 発火率は負にならない、$[u]_+$ & 第\ref{sec:glutcomputer}、\ref{sec:gaba}章 \\
オン--オフの対 & 整流器の対、2線式符号 & 一方は「増加」に、他方は「減少」に応答する & 第\ref{sec:glutcomputer}章、\eqref{eq:dualrail} \\
シャント抑制を受けるニューロン & アナログ除算器、ゲイン制御 & 入力から引くのではなく、入力を割る & 第\ref{sec:gaba}章、\eqref{eq:shunt} \\
\midrule
\multicolumn{4}{l}{\emph{発振器と時間}} \\
ペースメーカー & 弛張発振器、マルチバイブレータ & 一定の頻度で自らスパイクを出す & 第\ref{sec:serocomputer}、\ref{sec:dofa}章 \\
入力付きペースメーカー & 電圧制御発振器 & 固有のリズムを入力が速めたり遅めたりする & 第\ref{sec:chemoutput}章 \\
バースト発火ニューロン & ゲート付きバースト発生器 & 高頻度スパイクのバースト、つまり「ツー」を出す & 第\ref{sec:code}章、\eqref{eq:burst} \\
位相固定ニューロン & 位相検出器、位相同期ループ & 入力波の特定の位相で発火する & 第\ref{sec:sensors}、\ref{sec:chemoutput}章 \\
遅延のある軸索 & 遅延線 & 信号を決まった時間だけ遅らせる & 第\ref{sec:glutcomputer}章、図~\ref{fig:itd} \\
\midrule
\multicolumn{4}{l}{\emph{記憶と切り替え}} \\
双安定ニューロン & シュミットトリガ、ラッチ & 短い入力で長く続く状態に移る & 第\ref{sec:bistable}節 \\
樹状突起の枝をもつニューロン & マルチプレクサ、小さな多層ネットワーク & 許可入力があるときだけ枝が入力を通す & 第\ref{sec:synthesis}章 \\
\end{longtable}
\end{center}

表全体より重要な但し書きが一つある。これらは別々の細胞ではなく、別々の\emph{モード}である。同じニューロンが、入力が遅ければ積分器になり、抑制が$\tau$を短くすれば同時検出器になる（第\ref{sec:code}章）。覚醒中はペースメーカーで、睡眠中は黙った受信器になる。どの素子になるかは、膜のイオンチャネルと、それを調整するブロードキャストチャネルが決める。

\section{本書にまだ出ていない四つの素子}

\subsection{共振器}
\label{sec:resonator}

一部のニューロンには、電圧のあらゆる変化に逆らう遅い電流がある。電圧が上がれば電流はそれを引き下げ、下がれば引き上げる。ただし遅れを伴う。電子技術者から見ると、この電流はインダクタンスのように振る舞い、膜容量と合わせて共振回路を作る。ニューロンは特定の周波数、ふつう数Hzから十数Hzの入力に最も強く応答し、それより遅い入力にも速い入力にも弱く応答する~\cite{HutcheonYarom2000}。これは1個の細胞の中のバンドパスフィルタであり、ノイズの多い入力から自分の周波数のリズム、たとえば第\ref{sec:gaba}章の局所リズムを選び出す。

\subsection{漏れのない積分器}
\label{sec:perfectint}

漏れ積分器が入力を覚えているのは時間$\tau$、つまり数十ミリ秒だけである。しかし眼を静止させておくには、自分の位置を数秒間覚えている必要がある。「回転せよ」という指令は短時間の速度として来るが、眼は新しい位置に保たなければならない。ネットワークはこれを、第\ref{sec:glutcomputer}章の記憶と同じフィードバックで解く。時定数$\tau$のニューロン群が重み$w$で自らを興奮させるなら、$\tau\,\dd r/\dd t=-r+w\,r+I$であり、群の時定数は
\begin{empheq}[box=\keybox]{equation}
\tau_{\text{eff}} \;=\; \frac{\tau}{1-w} .
\label{eq:perfectint}
\end{empheq}
となる。$\tau=0.1\,\text{s}$、$w=0.995$なら$20\,\text{s}$になる。個々には10分の1秒で忘れるニューロンから作った、ほぼ理想的な積分器である~\cite{Seung1996}。代償は精密な調整である。$w$が1をわずかに超えれば積分器は飽和し、わずかに下回ればすぐに忘れる。そのため、このような積分器には第\ref{sec:motorlearning}章で述べたような学習による絶え間ない調整が必要である。

\subsection{双安定ニューロン}
\label{sec:bistable}

第\ref{sec:glutcomputer}章と第\ref{sec:gaba}章では、フリップフロップをニューロン群で組み立てた。しかし1個の細胞の中のフリップフロップもある。一部のニューロンには、いったんオンになると自ら興奮を維持する電流がある。短い入力でニューロンは長く続く発火状態、すなわち\emph{プラトー}に移り、抑制で静止状態に戻る。これはシュミットトリガである。ニューロンは二つの安定状態とその間のヒステリシスをもつ。たとえば筋肉を制御する\nt{ACET}ニューロンがこう振る舞い、この性質はブロードキャストチャネルでオンになる。セロトニンがニューロンに届くとプラトーが現れる~\cite{Hounsgaard1988}。同じ細胞が、\nt{SERO}があればラッチ、なければ普通の積分器である。

\subsection{積分器型と共振器型}

理論は興奮性細胞を二つのクラスに分ける~\cite{Izhikevich2007}。\emph{積分器型}は、ゼロから始まる電圧制御発振器に似ている。閾値をわずかに超えたニューロンはいくらでもゆっくり発火でき、発火率は入力とともに滑らかに上がる。このようなニューロンは届いたものをすべて足し合わせ、入力の大きさを発火率でよく伝える。\emph{共振器型}は、最低周波数のある発振器に似ている。ゆっくり発火することはできず、閾値の入力ですぐに有限の頻度でスパイクを出し、自分の周波数の入力を好む。前者は第\ref{sec:code}章のレート符号の自然な素子であり、後者はタイミング符号の自然な素子である。

\begin{keyblock}
\begin{proposition}[1個のニューロンは多くの素子になる]
\label{prop:eetypes}
電気的機能から見ると、ニューロンは漏れのある積分器と漏れのない積分器、同時検出器、ハイパスフィルタとバンドパスフィルタ、電圧--周波数変換器、整流器、除算器、発振器、遅延線、フリップフロップとして働く。ある細胞がどの素子になるかは、そのイオンチャネルと、それを調整するブロードキャストチャネルが決める。モード自体は測定済みである。脳がこうした切り替えを計算にどこまで使っているかは、主に仮説である。
\end{proposition}
\end{keyblock}

エンジニアにとって、これはプログラマブルロジックアレイに似ている。ハードウェアは一つで、素子の働きは構成で決まる。ただしここでは構成を書き込み時にではなく動作中に変える。変えるのは、第\ref{sec:serocomputer}--\ref{sec:acet}章で扱った遅いブロードキャストチャネルである。

\section{問題}

\begin{exercise}[\lvlA]
\label{ex:18:1}
\emph{漏れのない積分器の許容誤差。}$\tau=0.1$~sのニューロンが~\eqref{eq:perfectint}に従って眼の位置を$T=1$~sのあいだ、どちら向きにも$5\%$以内の誤差で保つ。(a)~$w$の許容範囲を求めよ。(b)~$w$は$K$個のシナプスの重みの和で、各々に独立な$10\%$の誤差がある。和のばらつきが許容範囲に収まるのは$K$がいくつのときか。
\end{exercise}

\begin{exercise}[\lvlB]
\label{ex:18:2}
\emph{回路としての共振器。}第\ref{sec:resonator}節の遅い電流を変数$W$で表す：$C\,\dd V/\dd t=-g_LV-g_wW+I$、$\tau_w\,\dd W/\dd t=V-W$。これが並列枝$R_w=1/g_w$、$L_w=\tau_w/g_w$をもつ$RC$回路であることを示し、$C/g_L=10\ms$、$\tau_w=100\ms$、$\alpha=g_w/g_L=4$で共振周波数と、そこでの$|Z|$の$|Z(0)|$に対する利得を求めよ。
\end{exercise}

\begin{exercise}[\lvlB]
\label{ex:18:3}
\emph{積分器はどう発火し始めるか。}(a)~ニューロン$\tau\,\dd V/\dd t=-V+\mu$、$\tau=10\ms$、閾値$\theta$、ゼロへのリセット。発火率が$1$~Hzになるのは$\mu/\theta-1$がいくつのときか。(b)~モデル$\tau\,\dd v/\dd t=v^2+\mu$（スパイクは$v$の$+\infty$への発散、リセットは$-\infty$へ）は$\mu=0.01$でどの発火率を与えるか。
\end{exercise}

\begin{exercise}[\lvlB]
\label{ex:18:4}
\emph{シミュレーション：積分器型と共振器型。}問題~\ref{ex:18:2}のモデルで$g_L=1$、$\tau_m=10\ms$、$\tau_w=100\ms$、閾値$1$、リセット$V\to0$（$W$は変えない）とし、入力$\mu=1.5\sin2\pi ft$、$f=1$--$40$~Hzを与える。$\alpha=0$と$\alpha=4$で、ニューロンがスパイクを出す周波数帯はどこか。条件$1.5\,|Z(f)|>1$と比べよ。
\end{exercise}

\begin{exercise}[\lvlB]
\label{ex:18:5}
\emph{1個の細胞のラッチ。}プラトー電流をもつニューロン：$\tau\,\dd r/\dd t=-r+F(wr+I)$、$F(x)=\min(1,\max(0,x))$、$\tau=10\ms$、$w=1.5$、背景入力$I=-0.2$。(a)~パルス$I=0.3$がニューロンを$r=0$からプラトーに移すのに必要な最短の長さはいくらか。(b)~長いパルス$I=-0.2-B$がニューロンを静止状態に戻すのに必要な最小の振幅$B$はいくらか。
\end{exercise}

\begin{exercise}[\lvlB]
\label{ex:18:6}
\emph{積分器の自己調整。}注視中は入力$I=0$で、滑りセンサがずれの速度$e=\dd r/\dd t$を与え、$\dd w/\dd t=-\eta\,e\,r$とする。$w=0.95$、$\tau=0.1$~s、$\langle r^2\rangle=1$のとき、$60$~sの注視で$|1-w|$が問題~\ref{ex:18:1}の許容範囲に入る$\eta$を求めよ。眼の回転中にも学習すると何が壊れるか。
\end{exercise}

\begin{exercise}[\lvlC]
\label{ex:18:7}
\emph{なぜ素子を組み替えるのか。}ブロードキャストチャネルが問題~\ref{ex:18:2}のモデルの$\alpha$を$0$と$4$の間で切り替える。電流の白色ノイズ中の弱い正弦波について、$11$~Hzと$1$~Hzで共振器型は積分器型よりSN比で何倍有利か。モードの切り替えそのものが有利になる課題を考えよ。
\end{exercise}
